\documentclass[a4paper,12pt]{article}
\usepackage[margin=1in]{geometry}
\usepackage{graphicx}
\usepackage{listings}
\usepackage{fancyhdr}
\usepackage{xcolor}
\usepackage{titling}

% Set Times New Roman font
\usepackage{fontspec}
\setmainfont{Times New Roman}

% Page setup
\pagestyle{fancy}
\fancyhf{}
\fancyhead[L]{Cryptography and Network Security(3161606)}
\fancyhead[R]{Enrollment No.}
\fancyfoot[C]{\thepage}
\renewcommand{\headrulewidth}{0.4pt}
\renewcommand{\footrulewidth}{0.4pt}

% Code listing style
\definecolor{codegreen}{rgb}{0,0.6,0}
\definecolor{codegray}{rgb}{0.5,0.5,0.5}
\definecolor{codepurple}{rgb}{0.58,0,0.82}
\definecolor{backcolour}{rgb}{0.95,0.95,0.92}

\lstdefinestyle{mystyle}{
    backgroundcolor=\color{backcolour},   
    commentstyle=\color{codegreen},
    keywordstyle=\color{magenta},
    numberstyle=\tiny\color{codegray},
    stringstyle=\color{codepurple},
    basicstyle=\ttfamily\footnotesize,
    breakatwhitespace=false,         
    breaklines=true,                 
    captionpos=b,                    
    keepspaces=true,                 
    numbers=left,                    
    numbersep=5pt,                  
    showspaces=false,                
    showstringspaces=false,
    showtabs=false,                  
    tabsize=2
}

\begin{document}

% Title
\begin{center}
\Large\textbf{Practical No.11}
\end{center}

% Aim
\section*{Aim}
\large\textbf{To implement and demonstrate the Diffie-Hellman Key Exchange algorithm for secure key exchange between two parties.}
\normalsize

\vspace{0.5cm}

% Theory
\section*{Theory}
The Diffie-Hellman key exchange is a method of securely exchanging cryptographic keys over a public channel. It allows two parties to establish a shared secret key that can be used for encrypted communication, without transmitting the key itself.

\vspace{0.3cm}
\textbf{Algorithm Steps:}
\begin{enumerate}
\item Both parties agree on public parameters: a prime number $p$ and a generator $g$
\item Each party generates a private key: Alice ($a$), Bob ($b$)
\item Each party computes their public key: $A = g^a \mod p$, $B = g^b \mod p$
\item Parties exchange public keys
\item Each party computes the shared secret: $S = B^a \mod p = A^b \mod p$
\item Both parties now have the same shared secret $S = g^{ab} \mod p$
\end{enumerate}

\vspace{0.3cm}
\textbf{Security:} The security relies on the discrete logarithm problem - it's computationally infeasible to derive $a$ from $A = g^a \mod p$ for large prime numbers.

% Code
\section*{Code}

\begin{verbatim}
/**
 * Diffie-Hellman Key Exchange - Practical 11
 * Test case: p=23, g=5, Alice private key=6, Bob private key=15, Expected shared secret=2
 */

// ============ DIFFIE-HELLMAN IMPLEMENTATION ============

function generatePrivateKey(p: number): number {
    // Generate random private key in range [2, p-2]
    return Math.floor(Math.random() * (p - 3)) + 2;
}

function computePublicKey(privateKey: number, g: number, p: number): number {
    // Compute public value: g^privateKey mod p
    return Math.pow(g, privateKey) % p;
}

function computeSharedSecret(publicKey: number, privateKey: number, p: number): number {
    // Compute shared secret: publicKey^privateKey mod p
    return Math.pow(publicKey, privateKey) % p;
}

function diffieHellmanKeyExchange(): number {
    // Public parameters (typically much larger in practice)
    const p = 23;  // Prime modulus
    const g = 5;   // Generator (primitive root)

    console.log("Public parameters:");
    console.log(`  Prime (p): ${p}`);
    console.log(`  Generator (g): ${g}`);

    // Alice's key generation
    const a = generatePrivateKey(p);
    const A = computePublicKey(a, g, p);
    console.log("\nAlice:");
    console.log(`  Private key (a): ${a}`);
    console.log(`  Public value (A = g^a mod p): ${A}`);

    // Bob's key generation
    const b = generatePrivateKey(p);
    const B = computePublicKey(b, g, p);
    console.log("\nBob:");
    console.log(`  Private key (b): ${b}`);
    console.log(`  Public value (B = g^b mod p): ${B}`);

    // Exchange public values and compute shared secret
    const S_alice = computeSharedSecret(B, a, p);
    const S_bob = computeSharedSecret(A, b, p);

    console.log("\nShared Secret Computation:");
    console.log(`  Alice computes: S = B^a mod p = ${B}^${a} mod ${p} = ${S_alice}`);
    console.log(`  Bob computes:   S = A^b mod p = ${A}^${b} mod ${p} = ${S_bob}`);
    console.log(`  Shared secrets match: ${S_alice === S_bob}`);

    return S_alice;
}

// Run the Diffie-Hellman key exchange
console.log("Diffie-Hellman Key Exchange Demonstration");
console.log("=========================================");
const sharedSecret = diffieHellmanKeyExchange();
console.log("\nResult:");
console.log(`  Shared secret: S = ${sharedSecret}`);
console.log(`  Verification: Both parties computed the same value ✓`);

// Security analysis
console.log("\nSecurity Analysis:");
console.log("  Eavesdropper sees: p = 23, g = 5, A = " + Math.pow(5, 6) % 23 + ", B = " + Math.pow(5, 15) % 23);
console.log("  To find S, attacker must solve:");
console.log("    - Find a from: 5^a mod 23 = " + Math.pow(5, 6) % 23 + "  (discrete logarithm)");
console.log("    - Find b from: 5^b mod 23 = " + Math.pow(5, 15) % 23 + " (discrete logarithm)");
console.log("  Then compute: S = " + Math.pow(5, 6) % 23 + "^b mod 23 = " + Math.pow(5, 15) % 23 + "^a mod 23 = " + sharedSecret);
\end{verbatim}

% Output
\section*{Output}

\begin{figure}[h]
\centering
\includegraphics[width=\textwidth]{practical11-output.png}
\caption{Output of Diffie-Hellman Key Exchange Implementation}
\end{figure}

\end{document}
